\chapter{misc}

\comm{definir poliedral: lineal a trozos con finitos trozos}

\begin{teo}
Sea $\alpha\subset\R^n$, con $n\geq 4$ un arco localmente poliedral salvo un conjunto numerable $X\subset \alpha$. Entonces, $\alpha$ es plana en $\R^n$.
\end{teo}
\begin{proof}
\comm{falta justificar pq el arco estaría compuesto x numerables segmentos}
Extendiendo cada segmento de $\alpha$ obtenemos una familia de recta. Agregamos a esta familia una recta que pase por cada punto de $X$, llamamos $\{r_i\}$ a la familia. Cada par de estas rectas determina un plano $H_{ij}$ de dimensión
\end{proof}